% ============================================================
%  Chapter 9: Quantum Mechanics
% ============================================================
\chapter{量子力学：从公设到预言}
\label{ch:quantum}

\section{概述}

R4（量子公设）和 R5（Born 规则）的结合产生了量子力学的全部预言能力。本章从 R4b+R4c 出发推导 Schr\"{o}dinger 方程，从 R4c 出发推导 Planck-Einstein 关系和不确定性原理。

% ════════════════════════════════════════════════════════
\section{Schr\"{o}dinger 方程}
% ════════════════════════════════════════════════════════

\begin{lawbox}[Schr\"{o}dinger 方程]
  $i\hbar\frac{\partial\psi}{\partial t} = -\frac{\hbar^2}{2m}\nabla^2\psi + V(\mathbf{r})\psi$
  \quad\textbf{前提}：R4b（幺正演化）, R4c（正则对易）, R1
\end{lawbox}

\noindent\textbf{变量}：波函数 $\psi(\mathbf{r},t) = \langle\mathbf{r}|\psi(t)\rangle$ [L$^{-3/2}$], 动量算符 $\hat{\mathbf{p}} = -i\hbar\nabla$ [M$\cdot$L$\cdot$T$^{-1}$],
位置算符 $\hat{\mathbf{x}}$ [L], 哈密顿算符 $\hat{H} = \hat{\mathbf{p}}^2/(2m) + V(\hat{\mathbf{r}})$ [M$\cdot$L$^2\cdot$T$^{-2}$]。

\begin{derivationbox}[推导]
  \textbf{Step 1} — R4b 抽象演化方程：$i\hbar\frac{d}{dt}|\psi(t)\rangle = \hat{H}|\psi(t)\rangle$
  \hfill\texttt{[N]} 孤立系统，非测量

  \textbf{Step 2} — 投影到位置表象：$\psi(\mathbf{r},t) \equiv \langle\mathbf{r}|\psi(t)\rangle$

  \textbf{Step 3} — R4c 确定动量算符：$\langle\mathbf{r}|\hat{\mathbf{p}}|\psi\rangle = -i\hbar\nabla\psi$
  \hfill\texttt{[N]} $[\hat{x}_i,\hat{p}_j]=i\hbar\delta_{ij}$ 唯一确定（Stone-von Neumann）

  \textbf{Step 4} — 经典 $H$ 算符化：$\hat{H} = \hat{\mathbf{p}}^2/(2m) + V(\hat{\mathbf{r}})$
  \hfill\texttt{[N]} R4d（算符-可观测量对应）

  \textbf{Step 5} — 代入得：$i\hbar\frac{\partial\psi}{\partial t} = -\frac{\hbar^2}{2m}\nabla^2\psi + V\psi$
\end{derivationbox}

% ════════════════════════════════════════════════════════
\section{Planck-Einstein 关系}
% ════════════════════════════════════════════════════════

\begin{lawbox}[Planck-Einstein 关系]
  $E = h\nu$（光子能量量子）
  \quad\textbf{前提}：R4c（正则对易）\quad\textbf{推导号}：\texttt{deriv.energy\_frequency}
\end{lawbox}

\begin{derivationbox}[推导：量子谐振子到 Planck-Einstein 关系]
  \textbf{Step 1} — 量子谐振子：$\hat{H} = \hat{p}^2/(2m) + \frac{1}{2}m\omega^2\hat{x}^2$

  \textbf{Step 2} — 升降算符：$a = \sqrt{\frac{m\omega}{2\hbar}}\hat{x} + i\sqrt{\frac{1}{2m\hbar\omega}}\hat{p}$
  \hfill\texttt{[N]} $[\hat{x},\hat{p}] = i\hbar$

  \textbf{Step 3} — $[a, a^\dagger] = 1$，$\hat{H} = \hbar\omega(a^\dagger a + \frac{1}{2})$

  \textbf{Step 4} — 能谱 $E_n = \hbar\omega(n + \frac{1}{2})$，$\Delta E = \hbar\omega = h\nu$。单量子激发：$E_\gamma = h\nu$
\end{derivationbox}

% ════════════════════════════════════════════════════════
\section{de Broglie 波长}
% ════════════════════════════════════════════════════════

\begin{lawbox}[de Broglie 波长]
  $\lambda = h/p$（物质波）
  \quad\textbf{前提}：$E = h\nu$（光子类比 + 实验验证）
\end{lawbox}

\noindent\textbf{变量}：波长 $\lambda$ [L], 动量 $p$ [M$\cdot$L$\cdot$T$^{-1}$]。

\textbf{推导}：光子 $E = h\nu$, $p = E/c = h\nu/c = h/\lambda \implies \lambda = h/p$。
de Broglie (1924) 假设此关系对所有物质粒子成立——Davisson-Germer (1927) 电子衍射实验证实。

% ════════════════════════════════════════════════════════
\section{不确定性原理}
% ════════════════════════════════════════════════════════

\begin{lawbox}[不确定性原理]
  $\sigma_x \cdot \sigma_p \ge \hbar/2$；$\Delta E \cdot \Delta t \ge \hbar/2$
  \quad\textbf{前提}：R4c, R5, R4d, R1
\end{lawbox}

\noindent\textbf{变量}：标准差 $\sigma_A^2 \equiv \langle(\hat{A}-\langle\hat{A}\rangle)^2\rangle$，期望值 $\langle\hat{A}\rangle = \langle\psi|\hat{A}|\psi\rangle$ (R5)。

\begin{derivationbox}[推导]
  \textbf{Step 1} — Robertson-Schr\"{o}dinger 不等式：
  $\sigma_A^2\cdot\sigma_B^2 \ge |\frac{1}{2}\langle[\hat{A},\hat{B}]\rangle|^2$
  \hfill\texttt{[N]} $\hat{A},\hat{B}$ 厄米 (R4d)；期望值由 R5 定义

  \textbf{Step 2} — 代入 $[\hat{x},\hat{p}] = i\hbar$ (R4c)：$\sigma_x\cdot\sigma_p \ge |\frac{1}{2}\langle i\hbar\rangle| = \hbar/2$

  \textbf{Step 3} — 能量-时间（Mandelstam-Tamm）：
  $\Delta t \equiv \sigma_A/|d\langle A\rangle/dt|$，$d\langle A\rangle/dt = (i/\hbar)\langle[\hat{H},A]\rangle \implies \Delta E\cdot\Delta t \ge \hbar/2$
\end{derivationbox}

\texttt{reversible: false}——不确定性可来自其他原因（经典波包 Fourier 变换限制）。等号成立条件：Gaussian 波包（最小不确定态 = 相干态）。
